El resultado de la Ecuación \ref{eq:contradiccion}, que se dio bajo el criterio de realidad (la medición de Alice no perturba físicamente a la partícula distante de Bob), nos dice que ambas componentes deberían corresponder a elementos de realidad física. Esto lleva a EPR a sugerir que la Mecánica Cuántica proporciona una descripción incompleta de la realidad. \cite{EPR} Fue así que la expectativa de EPR era que una teoría más detallada reemplazara la Mecánica Cuántica en el futuro, una teoría que produjera los mismos resultados que los obtenidos a través de un análisis estadístico de algunas "variables ocultas".
 
Lo importante del método que retomamos para deducir la Ecuación \ref{eq:des_bell} es que surge de las suposiciones de realismo local y de las anticorrelaciones perfectas del estado singlete sin necesidad de introducir una forma matemática específica para las variables ocultas. Por lo tanto, cualquier teoría de variables ocultas locales que satisfaga las condiciones que nosotros tomamos también debe cumplir la desigualdad de la Ecuación \ref{eq:des_bell}, independientemente de los detalles particulares de dicha teoría.

El resultado más importante aparece al comparar la desigualdad de Bell, Ecuación \ref{eq:des_bell}, con las correlaciones que predice la Mecánica Cuántica para el estado singlete, dadas por la Ecuación \ref{eq:correlacion_MC}. Como mostramos anteriormente, existen direcciones de medición para las cuales las predicciones cuánticas violan dicha desigualdad. Esto significa que las correlaciones cuánticas no pueden ser reproducidas por una teoría de variables ocultas locales que satisfaga las suposiciones utilizadas en nuestra deducción. \cite{Bell}

Esta contradicción es precisamente la importancia del resultado de Bell: mientras que EPR sugería que la Mecánica Cuántica era incompleta y que podría existir una descripción más detallada de la realidad, Bell demostró que una teoría de variables ocultas locales no puede reproducir todas las predicciones cuánticas. Por lo tanto, no es posible mantener simultáneamente las suposiciones consideradas y reproducir las correlaciones del estado singlete. Sin embargo, esto no implica que las partículas puedan utilizarse para transmitir información instantáneamente, sino que las correlaciones entre sus resultados no admiten una explicación mediante el modelo local considerado.


Finalmente, aunque nuestro desarrollo se centró en dos partículas de espín $\frac{1}{2}$, este formalismo también puede aplicarse al estudio del entrelazamiento entre sistemas atómicos. En particular, es posible considerar dos niveles internos de cada átomo como un sistema cuántico efectivo de dos estados, análogo al espín $\frac{1}{2}$. De esta manera, dos átomos preparados en un estado entrelazado adecuado pueden presentar correlaciones cuánticas semejantes a las estudiadas anteriormente. Así, mediante la elección de distintas bases de medición, las desigualdades de Bell permiten evaluar si dichas correlaciones pueden ser descritas mediante una teoría de variables ocultas locales.
